%% =============================================================
%%  chapters/03_related_work.tex
%% =============================================================
\chapter{Related Work and State of the Art}
\label{ch:related-work}

\section{Pressure-Based Occupant Detection and Classification}
\label{sec:related-pressure}
Early automotive seat occupancy detection relied on simple threshold-based
binary weight switches, which were later extended to spatially resolved
pressure mats capable of distinguishing occupant classes beyond a simple
occupied/unoccupied signal \cite{mueller2019pressure}. Subsequent work
applied regression neural networks directly to pressure heatmaps to
estimate continuous occupant weight, demonstrating that spatial pressure
distribution carries substantially more information than aggregate load
alone \cite{nakamura2021weight}. A recent survey consolidates these
pressure-sensing approaches and highlights classification accuracy,
sensor cost, and robustness to seat cover variation as the primary axes
along which such systems are compared \cite{fernandez2022occupant}.

\section{Camera-Based Cabin Monitoring}
\label{sec:related-camera}
Camera-based approaches to cabin monitoring have progressed from
frontal-face driver drowsiness detection toward full-cabin, multi-occupant
pose estimation. Infrared camera systems combined with real-time pose
estimation pipelines have been shown to operate reliably across day and
night conditions inside the vehicle cabin \cite{schmidt2020driver}.
Lightweight convolutional architectures tailored to the automotive
embedded compute budget have been proposed specifically for occupant pose
classification \cite{patel2021occupant}, while RGB-D fusion methods
combine colour and depth cues to improve robustness against variable
illumination and occluding objects such as blankets or bags placed on
seats \cite{kowalski2023rgbd}.

\section{Multi-Modal Sensor Fusion Approaches}
\label{sec:related-fusion}
Beyond the automotive cabin monitoring domain, multimodal deep fusion has
been studied extensively for tasks such as autonomous driving perception
and human activity recognition. A broad survey categorises fusion
architectures by fusion level and highlights feature-level fusion with
learned attention mechanisms as consistently outperforming naive
data-level concatenation \cite{chen2020multimodal}. Cross-modal attention
fusion between camera and depth sensors has further been shown to improve
robustness specifically for human activity and pose recognition tasks
\cite{oliveira2022crossmodal}, providing direct architectural inspiration
for the camera-pressure fusion module proposed in this thesis.



\section{Synthetic Data Generation and Sim2Real Transfer}
\label{sec:related-sim2real}
Given the cost and privacy constraints of collecting large real pressure
sensor datasets across diverse occupant populations, synthetic data
generation using simulated human bodies and physics-based seat contact
models has emerged as a practical complement to real data collection
\cite{ivanov2021synthetic}. Sim2Real studies quantify the remaining
statistical gap between simulated and real sensor readings and propose
domain adaptation or fine-tuning strategies to close it; a recent study
specifically targeting seat pressure heatmaps reports meaningful
performance retention when models trained on a mixture of synthetic and
real data are evaluated on real-only test sets \cite{lindqvist2023sim2real}.

\ref{clever2020bodiesrest3dhuman} provides a comprehensive review of human body simulation techniques for generating synthetic pressure heatmaps, including parametric body models, pose libraries, and physics-based contact simulations.

\section{Gap Analysis and Positioning of This Work}
\label{sec:gap-analysis}
While the works reviewed above individually address pressure-based
classification, camera-based pose estimation, generic multimodal fusion,
or synthetic pressure data generation, none jointly combine camera and
pressure heatmap fusion with a synthetic data generation and Sim2Real
validation pipeline for the specific multi-task set of occupancy,
passenger classification, weight, and pose estimation. Table
\ref{tab:gap-analysis} summarises this positioning relative to the
reviewed literature.

\begin{table}[H]
\centering
\caption{Comparison of related work against the scope of this thesis.}
\label{tab:gap-analysis}
\begin{tabular}{@{}lllll@{}}
\toprule
\textbf{Work} & \textbf{Modalities} & \textbf{Outputs Estimated} & \textbf{Synthetic Data} & \textbf{Real-Time} \\
\midrule
\cite{mueller2019pressure} & Pressure & Occupancy class & No & Yes \\
\cite{clever2020bodiesrest3dhuman} & Pressure (sim) & Occupancy class & Yes & No \\
\cite{nakamura2021weight}  & Pressure & Weight & No & No \\
\cite{schmidt2020driver}   & Camera (IR) & Pose & No & Yes \\
\cite{patel2021occupant}   & Camera (RGB) & Pose class & No & Yes \\
\cite{kowalski2023rgbd}    & Camera (RGB-D) & Pose & No & Partial \\
\cite{oliveira2022crossmodal} & Camera + Depth & Activity class & No & Partial \\
\cite{ivanov2021synthetic} & Pressure (sim) & Occupancy class & Yes & No \\
\cite{lindqvist2023sim2real} & Pressure & Statistical validation & Yes & No \\
\textbf{This thesis} & Camera + Pressure & Occupancy, class, weight, pose & Yes & Target \\
\bottomrule
\end{tabular}
\end{table}
%% TODO: Revisit this table once the final list of baseline comparisons
%% used in Chapter 8 is fixed, and add any additional works found during
%% the systematic literature review.
